\ifdefined\iscombined
	\lecturetitle{Modules}
\else
\documentclass{article}
\usepackage{listings}
\usepackage{amsthm}
\usepackage{comment}
\usepackage{amsmath}
\usepackage{amsfonts}
\usepackage{sectsty}
\usepackage{hyperref}
\usepackage{fancyhdr}
\usepackage[top=1in,bottom=1in,left=1in,right=1in]{geometry}
\usepackage{float}
\usepackage{graphicx}
\usepackage{tabularx}
\usepackage{mathtools}

\date{
	\vspace{-.75em}
	{\large \today}
}

\title{
	\vspace*{-1.5em}
	{\Huge Lecture 7} \\
	\vspace{-1em}
	\rule{\textwidth/4}{.01em} \\
	\vspace{-.25em}
	{\Large Modules}
	\vspace{-.75em}
}

\author{
	{\large Matthew Farah}
}

\hypersetup{
	colorlinks=true,
	linkcolor=blue,
	filecolor=magenta,
	urlcolor=blue,
	citecolor=blue,
	anchorcolor=blue
}

\fancyhead{}
\fancyhead[L]{\today}
\fancyhead[R]{Lecture 7 - Modules}
\sectionfont{\fontsize{12}{12}\bfseries}
\subsectionfont{\fontsize{10}{12}\normalfont}
\subsubsectionfont{\fontsize{10}{12}\normalfont}

\begin{document}

\maketitle
\pagestyle{fancy}

\fi

\emph{I was really tired this lecture so I didn't take many notes, sorry.} \\

\begin{itemize}
	\item Systems are composed of modules.
	\begin{itemize}
		\item Modules are intended to address a single concern.
		\begin{itemize}
			\item Modules act as an abstraction over how the system deals with a concern.
			\begin{itemize}
				\item In this sense, modules hide information about the system.
				\item A modules `secret' is that which a module hides.
			\end{itemize}
			\item A single concern may be shared across multiple business events.
			\item Modules have minimal knowledge of each other.
		\end{itemize}
	\end{itemize}
	\item There are three types of modules:
	\begin{enumerate}
		%TODO: Explain what each do more clearly.
		\item Behaviour hiding modules.
		\item Decision hiding modules.
		\item Interface hiding modules.
	\end{enumerate}
\end{itemize}

\ifdefined\iscombined
	\newpage
\else
	\end{document}
\fi
